% ============================================================
% Section 84: 布洛赫振荡与齐纳隧穿
% ============================================================
\section{固体理论：布洛赫振荡与齐纳隧穿}
\label{sec:bloch-oscillation}

\begin{modelbox}[题目：一维晶格中电子在电场下的运动]
考虑一电荷为 $-e$ 的电子在一维点阵中，能量 $E(k)=-2T\cos ka$，$a$ 为晶格常数。有一小的均匀电场 $\mathcal{E}$ 施加在平行于晶格的方向上。定性地描述电子在 $k$ 空间和真实空间的运动：
\begin{enumerate}
    \item 无散射情况下；
    \item 有散射情况下。
\end{enumerate}
电场``小''的意思是什么？当电场不再是小量时，实际晶体（具有多支能带）中会发生什么现象？
\end{modelbox}

\begin{tipbox}[考点快速分析]
\begin{itemize}
    \item \textbf{准经典运动}：$\hbar\,dk/dt=F=-e\mathcal{E}$，$k$ 空间匀速漂移
    \item \textbf{布洛赫振荡}：$k$ 空间周期运动，真实空间交流振荡
    \item \textbf{散射的影响}：弛豫时间 $\tau$ 决定能否完成振荡
    \item \textbf{齐纳隧穿}：强电场下电子穿越能隙跃迁到高能带
\end{itemize}
\end{tipbox}

\begin{derivationbox}[标准解题过程]
\textbf{Step 1: 无散射情况（理想晶体）}

\textbf{$k$ 空间运动}：准经典方程
\begin{equation}
\hbar\frac{dk}{dt}=-e\mathcal{E}
\quad\Longrightarrow\quad
k(t)=k(0)-\frac{e\mathcal{E}}{\hbar}t.
\end{equation}
波矢随时间线性变化，电子在 $k$ 空间以恒定速率沿电场反方向匀速移动。

当电子到达布里渊区边界 $k=\pi/a$ 时，由于能带的周期性 $E(k)=E(k+2\pi/a)$，该状态等价于 $k=-\pi/a$。电子被``反射''并在另一端重新出现，继续匀速运动。$k$ 空间运动是周期性的，\textbf{布洛赫振荡周期}
\begin{equation}
\boxed{T_B=\frac{2\pi/a}{e\mathcal{E}/\hbar}=\frac{2\pi\hbar}{e\mathcal{E}a}.}
\end{equation}

\textbf{真实空间运动}：群速度
\begin{equation}
v(k)=\frac{1}{\hbar}\frac{dE}{dk}=\frac{2Ta}{\hbar}\sin ka.
\end{equation}
由于 $k(t)$ 线性变化，$v(t)$ 呈正弦振荡。电子在真实空间的位置
\begin{equation}
x(t)=\int v(t)\,dt\propto-\cos\bigl(ka(t)\bigr),
\end{equation}
即电子在真实空间作周期性振荡（\textbf{布洛赫振荡}），振幅 $\Delta x\sim 2T/e\mathcal{E}$，频率 $f_B=1/T_B=e\mathcal{E}a/(2\pi\hbar)$。

\textbf{结论}：无散射时，直流电场导致 $k$ 空间周期性运动和真实空间交流振荡。

\textbf{Step 2: 有散射情况（实际晶体）}

实际晶体中电子受声子、杂质、缺陷等散射，弛豫时间为 $\tau$。

若 $\tau\ll T_B$（散射频繁），电子在两次散射之间只能移动很小的波矢
\begin{equation}
\Delta k=\frac{e\mathcal{E}\tau}{\hbar}\ll\frac{\pi}{a},
\end{equation}
远不足以到达布里渊区边界。散射使电子波矢不断重新分布，阻碍布洛赫振荡的形成。电子表现为服从\textbf{欧姆定律}的漂移运动。

\textbf{布洛赫振荡的实验观测}：仅在极纯净的半导体超晶格中，且低温下（$\tau$ 足够长），才能通过太赫兹辐射等手段观测到。

\textbf{Step 3: ``小电场''的物理意义}

``小电场''指电场不足以引起显著的\textbf{带间隧穿}（齐纳隧穿，Zener tunneling）。电子在 $k$ 空间运动时，若能隙较小，电子有一定概率隧穿到更高能带，而非被反射回同一能带。

避免显著带间隧穿的临界条件为
\begin{equation}
\boxed{e\mathcal{E}a\ll\frac{E_{\text{gap}}^2}{E_F},}
\end{equation}
其中 $E_{\text{gap}}$ 为布里渊区边界处的能隙，$E_F$ 为费米能（或带宽量级）。此条件确保电子始终停留在同一能带内，准经典近似和单带有效质量理论成立。

\textbf{Step 4: 强电场下的现象}

当电场不再是小量（$e\mathcal{E}a\sim E_{\text{gap}}^2/E_F$）时：
\begin{itemize}
    \item 带间隧穿概率显著增大，电子可跃迁到更高能带
    \item 发生\textbf{齐纳击穿}（Zener breakdown），晶体的导电性急剧变化
    \item 绝缘体可变为导体，这一现象在\textbf{齐纳二极管}中被实际应用
\end{itemize}
\end{derivationbox}

\begin{formulabox}[答案汇总]
\begin{center}
\begin{tabular}{lll}
\toprule
\textbf{情况} & \textbf{$k$ 空间运动} & \textbf{真实空间运动} \\
\midrule
无散射 & 匀速运动，布里渊区边界反射，周期 $T_B=2\pi\hbar/(e\mathcal{E}a)$ & 周期性交流振荡（布洛赫振荡） \\
有散射 & $\Delta k=e\mathcal{E}\tau/\hbar\ll\pi/a$，不形成振荡 & 漂移运动，服从欧姆定律 \\
\bottomrule
\end{tabular}
\end{center}

\textbf{``小电场''}：$e\mathcal{E}a\ll E_{\text{gap}}^2/E_F$，避免齐纳隧穿。

\textbf{强电场}：发生齐纳隧穿，导致齐纳击穿。
\end{formulabox}

\begin{insightbox}[物理图像]
\begin{itemize}
    \item 布洛赫振荡是周期势场中电子在外电场下的\textbf{本征量子力学响应}，但通常在普通金属和半导体中无法观测，因为电子在飞秒--皮秒时间尺度内就被散射，远未完成一个振荡周期。
    \item 超晶格（人工周期结构，晶格常数 $a$ 可达纳米量级）可降低振荡频率、延长振荡周期，使得在低温高纯样品中观测布洛赫振荡成为可能。
    \item 齐纳隧穿是理解绝缘体击穿、齐纳二极管稳压原理的关键机制。它本质上是外电场诱导的能带间量子隧穿，与单粒子隧穿（如扫描隧道显微镜）不同，是集体电子态的跃迁。
\end{itemize}
\end{insightbox}
